\uuid{qzr7}
\niveau{PCSI}
\module{Analyse}
\chapitre{Ensembles et applications}
\sousChapitre{Fonctions}
%!TeX root=../../../encours.nouveau.tex
%%% Début exercice %%%

\duree{10}
\difficulte{2}
\auteur{Antoine Crouzet}
\datecreate{2024-12-01}
\titre{Des ensembles en bijection}
\contenu{
\texte{Montrer que les ensembles suivants sont en bijection, en explicitant une telle bijection :}
\question{$\left[0,\,1\right]$ et $\left[a,\,b\right]$.}
\reponse{Pour chacun des cas, il faut d'une part déterminer une fonction puis montrer qu'elle est bijective.
Soit $a<b$. Prenons $f:\left[0,\,1\right] \to \left[a,\,b\right]$ définie par $f:x\mapsto (b-a)x+a$. Alors $f$ est strictement croissante (car $b-a>0$) donc injective. De plus, soit $y\in \left[a,\,b\right]$ :
	\begin{align*}
		(b-a)x+a = y &\Leftrightarrow (b-a)x = y-a \\
		&\Leftrightarrow x =\frac{y-a}{b-a} \in \left[0,\,1\right] \text{ car } y\in \left[a,\,b\right]
	\end{align*}}
\question{$\left]0,\,1\right]$ et $\left[0,\,+\infty\right[$.}
\reponse{Soit $f:\left]0,\,1\right] \to \left[0,\,+\infty\right[$ définie par $f:x\mapsto \frac{1}{x}-1$. $f$ est bijective; en effet, soit $y\in \left[0,\,+\infty\right[$ :
	\begin{align*}
		f(x)=y &\Leftrightarrow \frac{1}{x}-1=y\\
		&\Leftrightarrow \frac{1}{x}=y+1\\
		&\Leftrightarrow x=\frac{1}{y+1} \text{ car $y\geq 0$}
	\end{align*}
	On remarque que, si $y\in \left[0,\,+\infty\right[$, $y+1\in \left[1,\,+\infty\right[$ et donc $\frac{1}{y+1} \in \left]0,\,1\right]$. $f$ est bien bijective.}
\question{$\N$ et l'ensemble des entiers naturels pairs.}
\reponse{Prenons $f:\N \to \{2p,\,p\in \N\}$ définie par $f(n)=2n$. $f$ est bijective. Soit $y\in \{2p,\,p\in \N\}$; on écrit $y=2p$ pour $p\in \N$. Alors \[ f(n)=y \Leftrightarrow 2n=2p \Leftrightarrow n=p \in \N \]
	Ainsi, $f$ est bijective.}
\question{$\N$ et l'ensemble des entiers naturels impairs.}
\reponse{De la même manière, on introduit $f:\N \to \{2p+1,\,p\in \N\}$ définie par $f(n)=2n+1$.}
}
